\chapter{Conclusions and Future Work}

\section{Conclusion}


This dissertation developed optimized Nb$_3$Sn coatings on Cu substrates using the thermally evaporated Cu-Sn route. The following notable results were achieved:

\begin{enumerate}
    \item Cu-Sn First: $T_c = 16$~K with transition width $\Delta T_c < 1$~K, among the sharpest reported for Cu-compatible fabrication routes
    \item Nb First: Uniform morphology with $T_c = 15$~K and transition width $\Delta T_c = 2$~K
    %\item Thick Cu-Sn First: Onset $T_c = 17.7$~K, verifying that increased film thickness enables strain relaxation toward bulk $T_c$ for high CTE mismatch substrates (Cu).
    \item  Evidence suggesting that strain due to thermal contraction mismatch can be mitigated, e.g., by thick layers
\end{enumerate}

This dissertation also developed and demonstrated processes for optimizing Nb$_3$Sn coatings on copper substrates, including polishing, diffusion barriers, enrichment of the tin source, avoidance of unwanted diffusion reactions, and reaction parameters within the constraints imposed by the thin-film chamber (discussed below). 

The particular advantage explored by this dissertation is the use of thermal evaporation to adjust the Sn content of a deposited Cu-Sn layer flexibly. This allowed exploration of Cu-Sn compositions richer in Sn than standard bronze, which is the solid solution alloy of up to 15 wt.\% Sn in Cu. Sn compositions up to 46 at.\% were explored. The higher Sn content reproduced advantages seen in the manufacture of internal-tin Nb$_3$Sn superconducting wires, where higher tin content supported faster diffusion and higher tin activity in the growing Nb$_3$Sn film, leading to more of the film having composition at or close to the stoichiometric 25\% Sn. Since the superconducting properties are best at the stoichiometric Nb:Sn ratio, the ability to deposit Sn-rich Cu-Sn layers was found to be crucial to optimizing the superconducting properties of films.

While measurements of Nb$_3$Sn cavities in RF conditions are well established for high-gradient accelerator designs, hardly any measurements exist for Nb$_3$Sn films or cavities made on Cu or Cu-alloy substrates. This dissertation presents among the first measurements of RF performance of large-area Nb$_3$Sn films made on Cu substrates as a function of temperature and magnetic field. The results from both mushroom cavity and in-house RF hexagonal cavity experiments indicated that a cavity with Nb$_3$Sn-coated walls could attain a quality factor significantly higher than that of a cavity made from Cu. These measurements validated the excitement in the ADMX community over the potential use of Nb$_3$Sn superconducting cavities for dark-matter detectors.

Hexagonal cavity prototypes were designed, modeled, fabricated, coated with Nb$_3$Sn, and tested. The hexagonal prism geometry, adapted from previous work~\cite{Ahn:2021fgb, Braine:2024nzi}, accommodates RF testing in standard laboratory magnets and is compatible with planar thin-film deposition. RF modeling revealed that if hexagonal rather than cylindrical cavity geometries are used and the seams are located on the edges of the hexagon, the degradation of $Q$ from high-resistance seams can be reduced (see Fig.~\ref{fig:GapnoGapCylinderHexagon}).



%We designed and tested a hexagonal cavity prototype that operates within our laboratory's infrastructure, enabling $Q$ vs $B$ measurements up to 9~T. At 50~mK and zero field, our Nb$_3$Sn-coated cavity (Nb-first recipe) measured $Q_L = 77{,}000$ compared to $Q_L = 55{,}000$ for the bare Cu cavity, demonstrating 1.4$\times$ improvement. Accounting for RF leakage and residual Cu regions, the expected $Q$ for a fully Nb$_3$Sn-coated cavity becomes $Q_{\text{Nb}_3\text{Sn,corrected}} = 1.3 \times 10^5$ based on extrapolated measured results. Operation in a 9~T magnetic field was successfully demonstrated, though $Q$ degraded more rapidly with increasing field than expected. We attribute this to insufficient film thickness and surface preparation challenges rather than intrinsic material limitations.

%This work establishes thermal-evaporated Cu-Sn route Nb$_3$Sn as a viable Cu-substrate fabrication method using standard laboratory equipment. The demonstrated recipes show promise, and the cavity testing infrastructure provides a foundation for future optimization towards $Q>3\times10^5$ at $9$~T, which is required for increased scan rates in future axion detector experiments.

\section{Topics For Future Work}

This dissertation was completed under time and temperature constraints for the heat treatments available for the Nb$_3$Sn reaction. Reactions were carried out in the ultra-high vacuum environment of the thin-film sputtering chamber to avoid contamination. The negative impact of overheating the deposition chamber and burning out the chamber heater constrained the reaction temperature to $\sim750~^\circ$C and the reaction time to 3 hours. Therefore, the superconducting properties displayed by several films could have been improved by longer and/or hotter reactions. The use of a high-vacuum tube furnace was found to introduce contamination, suggesting that an ultra-high-vacuum tube furnace would be necessary for future work.

The sputter chamber used in this dissertation imposed geometric constraints on the profiles and curvatures of the pieces used to form the cavities. Future studies could benefit from access to deposition equipment capable of coating more complex surfaces, especially curved ones. Several approaches used by other research groups were reviewed in this dissertation, and a future apparatus to produce Nb$_3$Sn ADMX cavities could incorporate many of the best features of these approaches, including those presented here.

%You might have other things to say about future work with fridges, characterizations, ADMX, etc. IF you do, make them relevant to the conclusions in the dissertation.


\section{Challenges}

The wall non-uniformities are believed to result from poor cavity surface preparation prior to coating, and the thermal evaporators' lack of substrate rotation. The film near the interface between the two half-cavity pieces began to peel, suggesting that the substrate holder overlapped too much into the inner cavity surface (Fig.~\ref{fig:8pieceDelamination}). Future depositions will decrease the substrate holder mask size so the film can coat past the inner walls, reducing edge effects. With these improvements, this design can be improved with the current infrastructure. This design enables surface resistance and high-Q characterization in magnetic fields using standard laboratory equipment, allowing thin-film groups with limited resources to assess the quality of their films for high-field RF applications.